\begin{claim}{C.31}{Proved}{Track control and continuity at the fixed stratum}
Let $f_t$ be a based null-homotopic loop in
$\operatorname{Homeo}^+(\R^2)$, let $G$ be its generated group as in
C.28, and let $Y$ be the union of the nonopen $G$-orbits. This is a
closed subset of the plane. Define $Q_{s,t}$ on $Y$ by C.27 on each
intrinsic line orbit, by the descended formula of C.29 on each circle
orbit, and by the identity on point orbits. Then every $Q_{s,t}$ is
a bijection of $Y$ preserving each orbit, and
\[
 Q_{s,t}^{-1}=Q_{s,t}[f^{-1}],\qquad
 Q_{s,t}(x)\in\{f_u(x):0\leq u\leq1\},\qquad
 Q_{s,t}^{-1}(x)\in\{f_u^{-1}(x):0\leq u\leq1\}.
                                                        \tag{C.31.1}
\]
Here $Q[f^{-1}]$ means the same construction applied to the based
inverse loop $u\mapsto f_u^{-1}$, which generates the same group.
Both evaluation maps $(s,t,x)\mapsto Q_{s,t}^{\pm1}(x)$ are continuous
at every point of $[0,1]^2\times Z$, where
$Z=\{x:f_u(x)=x\text{ for all }u\}$ is the point-orbit stratum.
The construction commutes with every homeomorphism commuting with
all $f_u$. No continuity at a nonfixed orbit is asserted here.
\end{claim}
\begin{proof}
The union of open orbits is open, proving that $Y$ is closed.
C.28 supplies exactly the remaining point, circle and intrinsic-line
types. The circle restrictions have zero winding by C.29, since
the ambient loop is null-homotopic. Thus the formulas defining $Q$
are available on every orbit in $Y$, and are independent of the
chosen line or circle coordinates by their proved naturality.
Each formula gives a homeomorphism of its single intrinsic line
or circle, so together they give a bijection of $Y$.

We prove the inverse identity before making any continuity claim.
For increasing homeomorphisms of a line, pointwise inversion
reverses order and exchanges finite minimum and maximum:
\[
 (a\wedge b)^{-1}=a^{-1}\vee b^{-1},\qquad
 (a\vee b)^{-1}=a^{-1}\wedge b^{-1}.
                                                        \tag{C.31.2}
\]
For example, at a point $x$ the first inverse is the larger of
$a^{-1}(x)$ and $b^{-1}(x)$: at this larger point one of $a,b$
equals $x$, and the other is at least $x$.
The same argument works for compact families, whose extrema
and inverse extrema are attained. Therefore, with the notation
of C.27,
\[
 m_r^{-1}(x)=\max_{0\leq u\leq r} f_u^{-1}(x),\qquad
 M_r^{-1}(x)=\min_{0\leq u\leq r} f_u^{-1}(x).
                                                        \tag{C.31.3}
\]
The median of three line homeomorphisms is the lattice expression
$(a\wedge b)\vee(a\wedge c)\vee(b\wedge c)$; its dual expression
$(a\vee b)\wedge(a\vee c)\wedge(b\vee c)$ has the same value.
Applying (C.31.2) thus proves
\[
 \operatorname{med}(a,b,c)^{-1}
       =\operatorname{med}(a^{-1},b^{-1},c^{-1}).
                                                        \tag{C.31.4}
\]
Combining (C.31.3)--(C.31.4), and using symmetry of the median
in its two envelope entries, gives the first identity in (C.31.1)
on every line orbit. On a circle, the inverse of the unique based
lift $F_u$ is exactly the unique based lift of $f_u^{-1}$.
The same calculation on the universal line descends to the circle.
It also holds on point orbits.

At any point the median chooses one of its three entries. Each
envelope entry is $f_u(x)$ for a parameter attaining its extremum,
and the middle entry is $f_t(x)$ itself. This proves the second
assertion of (C.31.1) on lines; projection of the same statement
for the based lifts proves it on circles. The inverse identity
then proves the third assertion.

Fix $p\in Z$ and an open neighborhood $V$ of $p$ in the plane.
Joint continuity of the action, compactness of $[0,1]$, and
$f_u(p)=p$ give a neighborhood $W$ of $p$ such that
$f_u(W)\subseteq V$ for every $u$. The compact inverse family
gives the same assertion for all $f_u^{-1}$ after shrinking $W$.
By (C.31.1),
\[
 Q_{s,t}^{\pm1}(W\cap Y)\subseteq V\cap Y
                   \quad\text{for all }s,t\in[0,1].       \tag{C.31.5}
\]
Since $Q_{s,t}^{\pm1}(p)=p$, this proves joint continuity at every
$(s,t,p)$, uniformly in the two homotopy parameters.

A homeomorphism commuting with all $f_u$ permutes the $G$-orbits,
preserves their four types, and maps intrinsic lines by line
homeomorphisms by C.28. Naturality in C.27 and C.29 gives
$hQ_{s,t}h^{-1}=Q_{s,t}$ on all of $Y$. This proves the asserted
equivariance without any continuity assertion at the remaining
points.
\end{proof}

\paragraph{Scope of the gluing step.}
Formula (C.31.5) resolves continuity at the fixed stratum for the
orbitwise construction on $Y$. It does not yet make the bijections
$Q_{s,t}:Y\to Y$ homeomorphisms, because continuity at the nonfixed
line and circle strata remains unproved. Nor does it control a
contraction on the omitted open orbits. These are the remaining
global assembly problems for this route to C.12.
